% !PRINTED 704
\weberpage{704}
Daraus ergiebt sich, dass in dem Producte (6) die Theilproducte
\[
\prod \frac{1}{1-p^{-sf}},
\]
die sich über alle Primzahlen $p$ erstrecken, für die $f>1$ ist, auch
für $s=1$ noch endlich bleiben. Da andererseits $g, h$ bestimmte
von Null verschiedene Werthe haben, so ergiebt sich aus (7) der
wichtige Satz:

I. \gesperrt{Durchläuft $\mathfrak{p}$ die Gesammtheit der Primideale
ersten Grades irgend eines algebraischen Körpers $\Omega$, so hat das Product}
\weq{10}{(s-1)\prod\frac{1}{1-N(\mathfrak{p})^{-s}}}
\gesperrt{für $s=1$ einen endlichen von Null verschiedenen
Grenzwerth.}

Diese Eigenschaft bleibt auch dann noch erhalten, wenn bei
der Bildung des Productes eine beliebige \gesperrt{endliche} Anzahl von
Primidealen ersten Grades ausgelassen wird.

Die Normen $N(\mathfrak{p})$ sind in der Formel (10) gleich natürlichen
Primzahlen $p$. Eine Primzahl $p$ kommt aber darin so oft
vor, als sie Primfactoren ersten Grades in $\Omega$ enthält, also höchstens
$n$ mal. Ist $\Omega$ ein \gesperrt{Normalkörper}, so kommt auch wirklich
jede Primzahl $p$, die in Primfactoren ersten Grades zerlegbar
ist, genau $n$ mal darin vor, und wir können für das Product (10)
auch setzen:
\[
(s-1)\prod\frac{1}{(1-p^{-s})^n},
\]
worin sich das Product $\prod$ auf alle in \gesperrt{Primfactoren ersten
Grades zerlegbaren} Primzahlen $p$ erstreckt (\S.\,160).
% [source note] The assertion of exactly n distinct degree-one prime factors requires an unramified prime. For Q(i), 2 ramifies as a unit times (1+i)^2 and has only one distinct degree-one prime. Omitting the finitely many ramified Euler factors preserves the finite nonzero limiting-value conclusion, but the displayed products are not literally identical without that qualification.

Eine unmittelbare Folgerung des Satzes I. ist die:

\gesperrt{In jedem algebraischen Körper giebt es unendlich
viele Primideale ersten Grades.}

\begin{center}
\rule{6em}{0.4pt}
\end{center}
